\section{Background}
\label{sec:background-strands}

To interpret random forests it is common to calculate an variable importance measure in a similar vein to how a linear model regression estimates effect size. A commonly used measure of variable importance in random forest is the \emph{impurity importance}.
In this paper we introduce a series of experiments demonstrating how impurity importance differs from effect sizes in various regression settings. The motivation for this is that it is common (bad) practice interpret impurity performance as if it was a per-unit effect size. We demonstrate this through an LLM-based screening published articles that utilizes impurity performance to drive conclusions about effects without adequate additional corroboration.

In decision tree algorithms, a node's impurity measures the variability of the target variable among the samples reaching that node. 
For a random forest \citep{breiman2001}, the impurity importance of a feature is
the total decrease in node impurity attributable to splits on that feature,
summed over all splits and all trees \citep{breiman1984,louppe2013}. 
%
Formally, the Mean-Decrease-in-Impurity (MDI) for feature $j$ is defined as the decrease in node impurity across all splits on that feature, averaged over all trees:

\begin{equation}
\text{MDI}_j = \frac{1}{T}\sum_{t=1}^{T} \sum_{v \in S_j^{(t)}} p(v) \Delta i(v)
\end{equation}

\noindent
where $T$ is the number of trees, $S_j^{(t)}$ is the set of internal nodes in tree $t$ that split on feature $j$, $p(v)$ is the proportion of samples reaching node $v$, and $\Delta i(v)$ is the decrease in impurity (e.g., Gini or variance) caused by the split at node $v$.

This paper focus on effects of impurity importance in a regression setting.
For
regression, \code{ranger} \citep{wright2017}
(and most implementations) score a split by how much it reduces the sum of
within-child \textbf{response variances}. Note that 
for classification the impurity is the
Gini index, which measures the same thing: since $\mathbf{1}\{Y=k\}$ is
Bernoulli with variance $p_k(1-p_k)$,
$1 - \sum_k p_k^2 = \sum_k \mathrm{Var}(\mathbf{1}\{Y = k\})$, the total variance
of the class indicators. 
%Impurity importance is a variance-explained quantity in
%both settings, and that is the source of every effect studied here.
Hence, the underlying metric is variance in both regression and classification settings, impurity importance is fundamentally a variance-explained quantity. 

Two strands of prior work bear on this. The first documents that MDI is
\emph{biased} toward continuous and high-cardinality
features \citep{strobl2007,nembrini2018}; \citet{nembrini2018} give the
bias-corrected ``actual impurity reduction'' (AIR), exposed by \code{ranger} as
\code{impurity\_corrected}; further debiased MDI variants \citep{li2019,zhou2021}
address the same type bias. The second is theoretical: \citet{louppe2013}
decompose asymptotic MDI into per-variable contributions (for totally
randomised trees on categorical inputs), and \citet{scornet2023} proves that for
additive regression with independent uniform inputs, fitted with
\emph{truncated} (not fully grown) trees, MDI recovers a
\emph{variance-explained} decomposition, each
feature's score tracking $\mathrm{Var}(f_j(X_j))$, which for a linear term is
$\beta_j^2\,\mathrm{Var}(X_j)$. \citet{gromping2009} had already set
random-forest importances beside the $R^2$-decomposition importances of linear
regression, the comparison this paper makes throughout; and on the permutation
side, \citet{benard2022} show that mean decrease in accuracy is inconsistent
under dependence and give the Sobol-MDA, a consistent estimator of the total
Sobol index --- under independence, the same variance contribution normalised
by $\mathrm{Var}(y)$.

Our results are consequences of the \citet{scornet2023} identity, not new biases. 
Their value is
to show, among uncorrelated features that are all of one kind --- all continuous,
or all binary, so that a bias between feature types cannot separate them --- how
far a variance-contribution measure departs from a per-unit effect size, and that it
\emph{provably cannot} separate an effect's shape from a feature's distribution.
The new contribution is thus the sharp skew-versus-nonlinearity demonstration and
an estimate (Sections~\ref{sec:prevalence-method}--\ref{sec:prevalence-results})
of how often the misreading reaches published conclusions.

All experiments use \code{ranger} 0.18.0; the code that generates them and the
figures is available with the rest of the supporting material
(Section~\ref{sec:availability}). Unless noted, importances are the plain
\code{impurity} measure (the one the definition above describes);
Section~\ref{sec:linear} compares it against \code{impurity\_corrected}.
